
\subsection{Vertical weight from the loss of horizontal} \label{subsec:vert-from-hor-loss}
We can now use Theorem \ref{loss of hor weight} to show that the vertical weight at level $8p$ is a definite proportion of the total weight at level 0. 

We first observe
\begin{lem}\label{increase of the total}
 For any $n\in \N$,
we have:
$$
  \| W_\can^{\ver+\hor}(U \sm \KK)\|_1 \leq \| W_\can^{\ver+\hor}(U^n \sm \KK)\|_1+ C(p).
$$
\end{lem}
\begin{proof}
This follows immediately from Theorem \ref{thm:total-increase},
because $U^n \sm \KK \subset U \sm \KK$,
and the inclusion is homotopic to a homeomorphism. 
\end{proof}

Combining Theorem \ref{loss of hor weight} with Lemma \ref{increase of the total}, we immediately obtain:
\begin{lemma} \label{separation property}
 Let $f$ be a renormalizable quadratic-like map with period $p$.
Then there exists $M=M(p)$ such that
$$
     \| W^\ver_\can(U^{8p}\sm \KK)\|_1 \geq  \frac{1}{2}  \| W_\can^{\ver+\hor}(U \sm \KK) \|_1 ,
$$
provided $ \| W_\can^h (U\sm \KK)\|_1>M(p)$.
\end{lemma}

\subsection{A suitable version of the Covering Lemma} \label{subsec:covering-lemma}
We will now state (and prove) a version of the Covering Lemma\cite{covering lemma} in the language of the canonical weighted arc diagram. In the next subsection, we will use Theorem \ref{thm:covering} to ``push forward'' the $\wcan^\ver(U^{8p} \sm \KK)$ 
to get a lower bound on $\wcan^\ver(U \sm \KK)$ and thereby prove Theorem \ref{def vertical teaser}.

As Theorem \ref{thm:covering} is simply a restatement of the Quasi-Invariance Law of \cite{covering lemma}, the reader may wish to read the statement of this Theorem and then skip to the next section. 
\begin{theorem} \label{thm:covering}
Suppose that 
\begin{itemize}
\item
$U$ and $V$ are simply connected Riemann surfaces,
\item
$A = \bigcup_{i=1}^p A_i$ is a union of $p$ disjoint compact connected full subsets of $U$,
\item
$B = \bigcup_{i=1}^n B_i$ is the same on $V$,
and 
\item
$f\from U\to V$ is a holomorphic branched cover of degree $D$,
\end{itemize}
such that 
\begin{itemize}
\item
all the critical values of $f$ lie in $B$,
\item
for each $i$,
$A_i$ is a component of $f^{-1}B_i$,
and 
\item
the degree of $f|_{A_i}$ is $d$
(in the sense that $A_i$ and $B_i$ have open neighborhoods $A_i'$ and $B_i'$ such that $f\from A_i'\to B_i'$ is a degree $d$ branched cover).
\end{itemize}
Suppose further that 
\begin{equation} \label{eq:covering-compare}
\norm{\wcan^\ver(U \sm A)}_1 \ge \epsilon \norm{\wcan^{\ver + \hor}(V \sm B)}_1,
\end{equation}
and 
\begin{equation} \label{eq:covering-large}
\norm{\wcan^\ver(U \sm A)}_1 \ge M(\epsilon, D, p)
\end{equation}
Then 
\begin{equation} \label{eq:covering-conclude}
\norm{\wcan^\ver(U \sm A)}_1 \le \frac{12d^2}{\epsilon}\norm{\wcan^\ver{V \sm A}}_1.
\end{equation}
\end{theorem}

Before proving this theorem,
let us review the notation of \cite{covering lemma} and relate it to $\Wcan(U \sm A)$.
Accordingly, 
let $U$ and $A$ be as in the statement of Theorem \ref{thm:covering}.
We then define three numbers determined by $U$ and $A$:
$$
  X = \WW(U,\, \bigcup_{j=1}^N A_j);
$$

\begin{equation}\label{Y}
   Y = \sum_{j=1}^N \WW(U, A_j),
\end{equation}

$$
   Z = \sum_{j=1}^N  \WW ( U \sm \bigcup_{k \not= j } A_k ,\, A_j ).
$$

It is easy to see that $X\leq Y\leq Z$.
We further claim that
\begin{equation} \label{eq:X-bounds}
\wcan^\ver(U \sm A) \le X \le \wcan^\ver(U \sm A) + 4p
\end{equation}
and
\begin{equation} \label{eq:Z-bounds}
\wcan^{2\hor + \ver}(U \sm A) \le Z \le  \wcan^{2\hor + \ver}(U \sm A) + 4p^2
\end{equation}
To see the lower bound for $X$ in \eqref{eq:X-bounds},
we need only observe that every vertical leaf in the canonical lamination for $U \sm A$ contributes to the defining width for $X$,
so $X$ is at least the total width of this vertical lamination.
To see the upper bound,
we consider the extremal foliation $\FF_X$ for the defining path family for $X$.
Then the leaves of $\FF_X$ can lie in at most $2p$ homotopy classes $\alpha_i$,
and 
$$
W(\FF_X|_{\alpha_i}) \le \wcan(U \sm A,\alpha_i) + 2
$$
by maximality. 
The derivation of \eqref{eq:Z-bounds} is similar and left as an exercise for the reader. 

We can now reduce Theorem \ref{thm:covering} to the version of the Quasi-Invariance Law stated in Section 3.1.1 of \cite{covering lemma}.
\begin{proof}[Proof of Theorem \ref{thm:covering}]
We wish to think of the setting of our Theorem as an instance of the setting of the QI Law stated in Section 3 of \cite{covering lemma},
and specialized in Section 3.1.1. 
So we apply the latter,
letting $U$ and $V$ in \cite{covering lemma} be our $U$ and $V$,
letting $\Lambda_j' = \Lambda_j$ be our $A_j$,
and letting $B_j' = B_j$ be our $B_j$.
Then the $Y_U$ in Section 3.1.1 of \cite{covering lemma} is our $Y(U, A)$; we will use only that it is at least $X(U, A)$.
\emph{Because the critical values of $f$ lie in $B$},
the $Z_V = Z_V\{B_j, B_j', CV\}$ in Section 3.1.1 is simply our $Z(V, B)$,
and of course $X_V$ in \cite{covering lemma} is our $X(V, B)$. 

The context and direction of the two inequalities in Section 3.1.1 of \cite{covering lemma} are such that we can replace $Y_U$ with $X_U$ in them. 
We can then derive,
assuming $\norm{W^\ver(U\sm A)} \ge 4p^\epsilon$,
\begin{align*}
Z_V &\le \norm{\wcan^{2\hor+\ver}(V \sm B)}_1 + 4p^2 \quad&\text{by \eqref{eq:Z-bounds}} \\
&\le 2 \norm{\wcan^{\hor+\ver}(V \sm B)}_1 + 4p^2 \\
&\le \frac3\epsilon \norm{\wcan^\ver(U \sm A)}_1 &\text{by \eqref{eq:covering-compare}}\\
&\le \frac3\epsilon X_U &\text{by \eqref{eq:X-bounds}}. 
\end{align*}
This is the separation assumption (with $\xi :=  3/\epsilon$) in Section 3.1.1 of \cite{covering lemma};
we can then conclude,
assuming \eqref{eq:covering-large} (which implies $X_U \ge K(\xi, p, D)$),
$$
X_U \le \frac{6d^2}\epsilon X_V,
$$
which implies
$$
\norm{\wcan^\ver(U \sm A)}_1 \le \frac{6d^2}\epsilon(\norm{\wcan^\ver(V \sm B)}_1 + 4 p).
$$
Assuming then that $\wcan^\ver(U \sm A) \ge 48d^2p/\epsilon$,
we obtain \eqref{eq:covering-conclude}.
\end{proof}

\subsection{Vertical part has a definite weight} \label{covering-application}
We can now prove Theorem \ref{def vertical teaser}:
\begin{proof}[Proof of Theorem \ref{def vertical teaser}]
   If $\| W_\can^\ver (U\sm \KK)\|_1 \geq \| W_\can^\hor (U\sm \KK)\|_1 $,
there is nothing to prove. So, we assume that the opposite inequality holds,
and then
$$
   \| W_\can^\hor (U\sm \KK)\|_1 > \frac{1}{2} \| W_\can^{\ver+\hor} (U\sm \KK)\|_1.
$$
Then taking $M=M(p)$ as in Corollary \ref{separation property},
we conclude that
\begin{equation}\label{separation for f^n}
     \| W^\ver_\can(U^{8p}\sm \KK)\|_1 \geq  \frac{1}{2}  \| W_\can^{\ver+\hor}(U \sm \KK) \|_1 ,
\end{equation}
provided  $\| W_\can (U\sm \KK)\|_1>2M$.

We are now ready to apply Theorem \ref{thm:covering} 
to the map $f^{8p}\from U^{8p} \to U$,
where we let $A = B = \KK$.
We observe that all the topological hypotheses are satisfied, 
with $d = 2^{8}$ and $D = 2^{8p}$, 
and also \eqref{eq:covering-compare} holds with $\epsilon = 1/2$.
Moreover,
taking $\normone{\wcan^{\hor+\ver}(U \sm \KK)}$ sufficiently large and applying Lemma \ref{increase of the total},
we have 
$$
\normone{\wcan^{\hor+\ver}(U^{8p} \sm \KK)} \ge \normone{\wcan^{\hor+\ver}(U \sm \KK)} - C(p) \ge M(\epsilon, D, p),
$$
as required in \eqref{eq:covering-large}.
We then obtain,
as our form of \eqref{eq:covering-conclude},
\begin{align*}
\normone{\wcan^{\ver}(U \sm \KK)} &\ge (24 \cdot 2^{16})^{-1} \normone{\wcan^{\ver}(U^{8p} \sm \KK)}.
\end{align*}
The Lemma then follows from this and Lemma \ref{separation property}.
\end{proof}




